\documentclass[11pt, oneside]{article} 
\usepackage{geometry}
\geometry{letterpaper} 
\usepackage{graphicx}
	
\usepackage{amssymb}
\usepackage{amsmath}
\usepackage{parskip}
\usepackage{color}
\usepackage{hyperref}

\graphicspath{{/Users/telliott/Dropbox/Github-Math/figures/}}
% \begin{center} \includegraphics [scale=0.4] {gauss3.png} \end{center}

\title{Euclid's Elements}
\date{}

\begin{document}
\maketitle
\Large

%[my-super-duper-separator]

\subsection*{excircles}

In the figure below we see $\triangle ABC$ and its incircle on center $I$.  Three more circles or parts of them are drawn.  They have the property that each one is tangent to a different side of the original $\triangle ABC$ and also tangent to extensions of the other sides.

\begin{center} \includegraphics [scale=0.4] {excircle_intro1.png} \end{center}

In the next figure we draw just one excircle for $\triangle ABC$, the one tangent to side $a$ opposite $A$.

The easiest method to find the center of the excircle, called $I_a$, is to bisect the external angles at $B$ and $C$ and find where they intersect, then draw $AI_a$ 

\begin{center} \includegraphics [scale=0.4] {excircle_intro3.png} \end{center}

$I_a$ lies on the bisector of $A$.  (\emph{Proof}.  Drop perpendiculars to the extensions of the sides at $D$ and $E$, forming two congruent triangles by HL.  $\square$)

We are supposing that there is a single point where the excircle is tangent to side $a$ of $\triangle ABC$, the point $F$.  We must prove that $F$ exists and is tangent to the excircle.  

Choose $F$ such that $CE = CF$.  Draw $I_aF$.

Since we used the angle bisectors to find $I_a$ we have $\triangle I_aCE \cong I_aCF$ by SAS.

$I_a E \perp CE$ so $I_aF \perp BC$.

Also, $I_aF = I_aE$, so $I_aF$ is then a radius of the circle.

Since $IaF$ is perpendicular to $BC$ at $F$, $BFC$ is tangent, with $CE = CF$ (as we're given) and also $BD = BF$.

One neat property of the excircle can be seen more easily with an obtuse angle at $B$ or $C$.  We choose $B$.

\begin{center} \includegraphics [scale=0.35] {excircle_intro4.png} \end{center}

Find where $BC$ crosses the bisector at $S$.  Find the tangent from $S$ to the incircle at $T$.  $SX$ is also tangent to the incircle from $S$, so $SX = ST$.

Drop the perpendiculars from $I$ to the sides at $X$, $Y$, $Z$ and $T$.
 
Finally, extend $ST$ to meet side $AC$ at $B'$ and go back to meet side $AB$ at $C'$.

The labels give a hint of where we're going.

From $ST=SX$ we get congruent right triangles $\triangle SIT \cong \triangle SIX$ by HL.  As a result we have a total of four equal angles at $S$ and then two other angles equal, also by vertical angles.

There are two more congruent triangles $\triangle ABS \cong \triangle AB'S$ by ASA.

This has a number of consequences, but we will leave this here for you to work with as an exercise.  We'll come back to it in the next book.

$\square$

\end{document}